\section{Binary distance separation and coordinate fields}

Write \(d_P(p)=|\{|p-q|:q\in P\setminus\{p\}\}|\) and \(D(P)=\max_{p\in P}d_P(p)\).
A binary tree on \(P\) is separating if, at every split, any pin in that node has distinct distances to targets in different children, with the pin itself omitted.

\begin{theorem}[982-D01]
Let \(P\) consist of \(n\ge2\) distinct planar points with no three collinear. If \(P\) admits a separating binary tree, then \(D(P)\ge\lceil n/2\rceil\).
\end{theorem}

\begin{proof}
In a node \(S\) of size \(m\), let the two children have sizes \(b\le c=m-b\).
Count the distinct distances separately for each pin in \(S\) and each child. A target pair has at most two witnesses in \(S\), since its bisector is a line. For the \(b\)-child use \(r-1\le\binom r2\) on each nonempty distance class of size \(r\). The total contribution is at least \(b(m-1)-b(b-1)=bc\). For the \(c\)-child use \(2r-3\le\binom r2\), giving at least \((2c(m-1)-c(c-1))/3=c(m+b-1)/3\). Thus
\[
 3\sum_{p\in S}d_S(p)\ge(m-b)(m+4b-1).
\]
On the path to \(S\), let the discarded siblings have sizes \(t_j\) and set \(V=\sum_j t_j^2\). For pins in \(S\), each sibling contributes at least \(mt_j-t_j(t_j-1)\) by the same pair count. Distances in distinct sibling groups and inside \(S\) are disjoint at each pin by separation at the first split between the targets. Hence
\[
 3mD(P)\ge(m-b)(m+4b-1)+3(m+1)(n-m)-3V.
\]
At the root \(m=n\) and \(V=0\). Maintain \(m>n/2\) and \(3V\le m(n-m)\). If \(b<(2m-n)/4\), descend into the larger child. The new size \(m-b\) exceeds \(n/2\) and the slack in the second invariant increases by \(b(2m-n-4b)>0\). The positive integer size decreases, and a leaf cannot meet \(m>n/2\), so eventually \(b\ge(2m-n)/4\). At this stopping node
\[
 3m\bigl(D(P)+1/3-n/2\bigr)
 \ge m(n/2-m)+3mb-4b^2+b+3(n-m).
\]
The first three terms on the right are concave in \(b\), with values \((m-n/2)(n-m)/2\) at \((2m-n)/4\) and \(m(n-m)/2\) at \(m/2\). The right side is nonnegative. Since \(D(P)\) is an integer, \(D(P)\ge n/2-1/3\) implies the stated bound in both parities.
\end{proof}

\begin{theorem}[982-D02]
Let \(P\) be as in the preceding theorem. If a similarity takes \(P\) into \(\mathbb Q^2\), then \(D(P)\ge\lceil n/2\rceil\).
\end{theorem}

\begin{proof}
After a further scaling the coordinates are integers. Split a node according to the parity of \(x+y\). Modulo two the squared distance is the sum of the two points' parities, so different target classes give different distances from any pin. Within either child, translate an integer point of the child to the origin and apply
\((x,y)\mapsto((x+y)/2,(y-x)/2)\).
The output is integral and squared distances are halved. Apply the transformation also if the split has only one nonempty child. Every surviving pair has a positive integer squared distance, so no branch can halve indefinitely. Suppressing one-child nodes yields the tree required by 982-D01.
\end{proof}

\begin{proposition}[982-D03, 982-D04]
Let \(K\subset\mathbb R\) be a number field. If the quadratic extension \(K(i)/K\) ramifies at some finite place above \(2\), then every \(P\subset K^2\) as in 982-D01 satisfies \(D(P)\ge\lceil n/2\rceil\). In particular this holds for every odd-degree real number field, and the conclusion is invariant under a common similarity.
\end{proposition}

\begin{proof}
At a ramified place the local quadratic extension \(L=K_v(i)/K_v\) is totally ramified. Its residue field has characteristic two, and conjugation \(\sigma:i\mapsto-i\) acts trivially there. Identify \((x,y)\) with \(x+iy\); the squared distance for differences is the norm \(N(z)=z\sigma(z)\).
For each nonsingleton group translate by a member and divide by a power of a uniformizer of \(L\) chosen so that the minimum valuation among all nonzero differences becomes zero. This multiplies all norms of differences by the same nonzero element of \(K_v\). Reduction now gives at least two classes. If the pin has residue \(a\), then the norms of differences to targets with residues \(b\) and \(c\) reduce to \((a-b)^2\) and \((a-c)^2\). These differ for \(b\ne c\). Binary regrouping of distinct residue classes preserves separation. Each original residue child is strictly smaller. This proves the first claim.

At a place where \(v(2)\) is odd, the extension of the normalized valuation satisfies \(v(1+i)=v(2)/2\), since \((1+i)^2=2i\) and \(i\) is a unit. The nonintegral value forces ramification of degree two. In an odd-degree number field \(\sum_{\mathfrak p\mid2} e_{\mathfrak p} f_{\mathfrak p}=[K:\mathbb Q]\), so an odd \(e_{\mathfrak p}\) exists.
\end{proof}

\begin{corollary}[982-D05, 982-D06]
If a strictly convex real counterexample exists at size \(n\), one exists at that size with real algebraic coordinates. Every real number field containing the coordinates of a similar copy of any counterexample has even degree, all its ramification indices above \(2\) are even, and \(K(i)/K\) is unramified at these places.
\end{corollary}

\begin{proof}
For a counterexample, prescribe the signs of all oriented triple determinants and of all within-pin differences of squared distances. The prescriptions are a finite existential formula over \(\mathbb Q\). The transfer theorem for real closed fields gives a solution over the real algebraic numbers, preserving the supporting edges and all within-pin distance partitions. The local criterion and the valuation calculation above give the restrictions on every field containing the coordinates.
\end{proof}
